\documentclass[10pt,letterpaper]{article}
\usepackage{cornell-notes}

\title{Chapter 11: Case Study 2: Windows 8}
\author{}
\date{}

\setDocTitle{Chapter 11: Case Study 2: Windows 8}
\setDocSubtitle{Operating Systems}
\setDocVersion{v1.0}
\setDocStatus{Study Companion}
\setDocOwner{Operating Systems Cornell Notes}
\setDocAuthor{}
\setDocDate{\today}
\setDocSummary{Cornell-format study and review notes for Chapter 11: Case Study 2: Windows 8.}

\setCornellCollection{Operating Systems Cornell Notes}
\setCornellUnitType{Chapter}
\setCornellUnitNumber{11}
\setCornellUnitTitle{Case Study 2: Windows 8}
\setCornellCompanionLabel{Study and review companion}

\hypersetup{pdftitle={Chapter 11: Case Study 2: Windows 8},pdfsubject={Cornell notes},pdfauthor={}}

\begin{document}
\maketitle
\section*{Learning Objectives}
\begin{studybox}{By the end of this chapter, you should be able to}
\begin{enumerate}
  \item Trace the historical transition from MS-DOS-based systems to the NT architecture described through Windows 8.1.
  \item Distinguish native NT services, Win32 APIs, and registry configuration.
  \item Explain HAL, kernel, executive, drivers, object manager, subsystems, DLLs, and user-mode services.
  \item Trace job, process, thread, and fiber management and priority-based scheduling.
  \item Explain sections, working sets, page lists, demand paging, and caching.
  \item Trace asynchronous I/O through handles, I/O request packets, device stacks, and completion.
  \item Explain NTFS master-file-table records, attributes, allocation, directories, and journaling.
  \item Analyze source-era power management, access tokens, security descriptors, ACLs, integrity levels, and exploit mitigations.
\end{enumerate}
\end{studybox}

\begin{studybox}{Section Map}
\hyperlink{sec-11-1}{11.1} \quad \hyperlink{sec-11-2}{11.2} \quad \hyperlink{sec-11-3}{11.3} \quad \hyperlink{sec-11-4}{11.4} \quad \hyperlink{sec-11-5}{11.5} \quad \hyperlink{sec-11-6}{11.6} \quad \hyperlink{sec-11-7}{11.7} \quad \hyperlink{sec-11-8}{11.8} \quad \hyperlink{sec-11-9}{11.9} \quad \hyperlink{sec-11-10}{11.10} \quad \hyperlink{sec-11-11}{11.11}
\end{studybox}

\section*{Cornell Cue-and-Notes Area}
\small
\begin{longtable}{C N}
\rowcolor{navy}
\color{white}\textbf{Cue / Recall Prompt} &
\color{white}\textbf{Notes}\\
\endfirsthead
\rowcolor{navy}
\color{white}\textbf{Cue / Recall Prompt} &
\color{white}\textbf{Notes (continued)}\\
\endhead
\rowcolor{ice}\multicolumn{2}{l}{\hypertarget{sec-11-1}{}\textbf{Section 11.1: History of Windows Through Windows 8.1}}\\*\hline
\textbf{11.1.1 MS-DOS} & \textcolor{legacy}{\textbf{Historical or legacy context:}} MS-DOS provided a small single-user environment for early personal computers with limited protection, memory, and multitasking support.\\\hline
\textbf{11.1.2 DOS-based Windows} & \textcolor{legacy}{\textbf{Historical or legacy context:}} Graphical Windows releases layered cooperative and later preemptive facilities over DOS heritage, accumulating compatibility constraints.\\\hline
\textbf{11.1.3 NT-based Windows} & \textcolor{legacy}{\textbf{Historical or legacy context:}} The NT line introduced protected virtual memory, preemptive threads, security, portability layers, and an object-based executive, later becoming the common base.\\\hline
\textbf{11.1.4 Windows Vista} & \textcolor{legacy}{\textbf{Historical or legacy context:}} Vista revised driver, graphics, security, service, and user-account behavior while retaining the NT architectural foundation.\\\hline
\textbf{11.1.5 Modern Windows} & \textcolor{legacy}{\textbf{Historical or legacy context:}} The source-era Windows 8 and 8.1 material emphasizes touch-oriented interfaces, application packaging, energy awareness, and continuity with desktop Win32 applications.\\\hline
\rowcolor{ice}\multicolumn{2}{l}{\hypertarget{sec-11-2}{}\textbf{Section 11.2: Programming Windows}}\\*\hline
\textbf{API layers} & Applications usually call documented environment APIs. User-mode libraries translate many requests into native services, local computation, or calls to subsystem processes.\\\hline
\textbf{11.2.1 Native NT API} & \textcolor{legacy}{\textbf{Historical or legacy context:}} Native services expose object, process, memory, file, and synchronization operations close to the executive and are commonly reached through system libraries.\\\hline
\textbf{11.2.2 Win32 API} & Win32 presents handles, processes, threads, files, windows, registry, and synchronization procedures. A procedure is not necessarily a one-to-one kernel system call.\\\hline
\textbf{11.2.3 Registry} & The registry stores hierarchical configuration hives, keys, and values. Centralization supports common access and policy but makes consistency, backup, permissions, and corruption consequential.\\\hline
\rowcolor{ice}\multicolumn{2}{l}{\hypertarget{sec-11-3}{}\textbf{Section 11.3: System Structure}}\\*\hline
\textbf{11.3.1 Architecture} & \textcolor{legacy}{\textbf{Historical or legacy context:}} Kernel mode contains the hardware abstraction layer, low-level kernel, executive managers, and drivers. User mode contains applications, environment components, service processes, and libraries.\\\hline
\textbf{11.3.2 Booting} & Firmware and the boot manager select and load the system; the loader brings in the kernel, HAL, boot drivers, and registry state; initialization creates system processes and starts services.\\\hline
\textbf{11.3.3 Object manager} & Executive objects have types, names when placed in a namespace, reference counts, security descriptors, and operations. Per-process handle tables give validated references rather than exposing kernel addresses.\\\hline
\textbf{11.3.4 Subsystems and services} & DLLs provide application APIs; environment and service processes perform user-mode policy and services; local procedure calls and system calls cross boundaries.\\\hline
\rowcolor{ice}\multicolumn{2}{l}{\hypertarget{sec-11-4}{}\textbf{Section 11.4: Processes and Threads in Windows}}\\*\hline
\textbf{11.4.1 Concepts} & A job groups processes for limits and accounting; a process owns an address space and handles; a thread is the scheduled entity; a fiber is cooperatively scheduled by code within a thread.\\\hline
\textbf{11.4.2 Management APIs} & Creation APIs establish executable image, address space, initial thread, inherited handles, environment, and security context. Waiting and synchronization use object handles.\\\hline
\textbf{11.4.3 Implementation} & Executive process and thread objects combine with lower-level scheduling structures. The scheduler selects the highest-priority ready thread and uses priority changes, quanta, affinity, and per-processor queues.\\\hline
\rowcolor{ice}\multicolumn{2}{l}{\hypertarget{sec-11-5}{}\textbf{Section 11.5: Memory Management}}\\*\hline
\textbf{11.5.1 Concepts} & Processes reserve address ranges and commit storage; section objects represent mapped files, images, or shared memory; demand paging populates pages when accessed.\\\hline
\textbf{11.5.2 Memory APIs} & APIs allocate, free, protect, lock, map, unmap, inspect, and share virtual memory. File-mapping objects connect section semantics to application handles.\\\hline
\textbf{11.5.3 Implementation} & \textcolor{legacy}{\textbf{Historical or legacy context:}} Working sets bound resident pages; the manager trims sets and moves frames among free, zeroed, standby, modified, and bad lists. Page faults recover data from files, page storage, or zero-fill.\\\hline
\rowcolor{ice}\multicolumn{2}{l}{\hypertarget{sec-11-6}{}\textbf{Section 11.6: Caching in Windows}}\\*\hline
\textbf{Unified cached views} & The cache manager maps file data into kernel virtual addresses and cooperates with the memory manager for demand paging, read-ahead, lazy writing, and memory-pressure reclamation.\\\hline
\rowcolor{ice}\multicolumn{2}{l}{\hypertarget{sec-11-7}{}\textbf{Section 11.7: Input/Output in Windows}}\\*\hline
\textbf{11.7.1 Concepts} & I/O is object- and packet-based, asynchronous by design, and routed through stacks of device objects and drivers. File-system volumes also appear within the device model.\\\hline
\textbf{11.7.2 I/O APIs} & Applications open handles, issue synchronous or overlapped operations, wait on events or completion ports, control devices, and close references. User-level completion style does not change the asynchronous kernel path.\\\hline
\textbf{11.7.3 Implementation} & An I/O manager creates an I/O request packet and sends it down a device stack. Drivers may complete, queue, transform, or forward it; interrupts and deferred procedures advance hardware work before completion unwinds upward.\\\hline
\rowcolor{ice}\multicolumn{2}{l}{\hypertarget{sec-11-8}{}\textbf{Section 11.8: The Windows NT File System}}\\*\hline
\textbf{11.8.1 Concepts} & \textcolor{legacy}{\textbf{Historical or legacy context:}} NTFS supports large files and volumes, Unicode names, multiple data streams, compression, encryption, access control, reparse points, and journaling.\\\hline
\textbf{11.8.2 Implementation} & Every file and directory has a master file table record containing typed attributes. Small attributes may be resident; large data use runs of clusters. Directories use indexed structures, and a log protects metadata transactions.\\\hline
\rowcolor{ice}\multicolumn{2}{l}{\hypertarget{sec-11-9}{}\textbf{Section 11.9: Windows Power Management}}\\*\hline
\textbf{Coordinated power states} & \textcolor{legacy}{\textbf{Historical or legacy context:}} The power manager coordinates system and device states through policy, idleness, driver requests, sleep or hibernation transitions, and device-stack power packets.\\\hline
\rowcolor{ice}\multicolumn{2}{l}{\hypertarget{sec-11-10}{}\textbf{Section 11.10: Security in Windows 8}}\\*\hline
\textbf{11.10.1 Concepts} & \textcolor{legacy}{\textbf{Historical or legacy context:}} Logon creates an access token containing identities, groups, privileges, and integrity information. Objects carry security descriptors with owner and access-control data.\\\hline
\textbf{11.10.2 Security APIs} & APIs authenticate, impersonate, enable privileges, inspect or set security descriptors, manipulate ACLs, and create restricted tokens, subject to the caller's rights.\\\hline
\textbf{11.10.3 Implementation} & The security reference monitor compares a token with an object's discretionary ACL, mandatory integrity policy, requested access, and privileges; auditing rules can record decisions.\\\hline
\textbf{11.10.4 Mitigations} & \textcolor{legacy}{\textbf{Historical or legacy context:}} Address-space randomization, data-execution prevention, stack protection, safer service isolation, volume encryption, signing, and restricted application environments raise exploitation cost.\\\hline
\rowcolor{ice}\multicolumn{2}{l}{\hypertarget{sec-11-11}{}\textbf{Section 11.11: Summary}}\\*\hline
\textbf{Integrated case study} & The NT architecture combines a hardware abstraction layer, kernel scheduler, object-based executive, layered drivers, demand-paged memory, mapped-file cache, NTFS metadata, and token-based security.\\\hline
\end{longtable}
\normalsize

\section*{Chapter Synthesis}
\begin{studybox}{Chapter Summary}
The source-era Windows case study organizes operating-system services around executive objects and handles. Threads are scheduled by priority, sections unify mapped images and shared memory, working sets and page lists govern residency, IRPs traverse device stacks, NTFS represents files as attribute collections in the MFT, and access tokens are checked against object security descriptors.
\end{studybox}

\begin{studybox}{Key Relationships}
\begin{itemize}
  \item A handle indexes a process table entry that references an executive object and carries granted access.
  \item Section objects connect executable images, mapped files, shared memory, the memory manager, and the cache manager.
  \item I/O request packets connect application operations to layered drivers and asynchronous hardware completion.
  \item Access tokens represent subject authority, while security descriptors represent object protection policy.
\end{itemize}
\end{studybox}

\begin{studybox}{Design Tradeoffs}
\begin{itemize}
  \item Centralized registry configuration simplifies discovery and policy but creates a sensitive shared dependency.
  \item Priority scheduling improves urgent response but requires boosting or policy controls to prevent starvation and inversion.
  \item Lazy cached writes improve throughput but require careful flush, ordering, and recovery behavior.
  \item Rich NTFS features improve functionality while increasing metadata and implementation complexity.
\end{itemize}
\end{studybox}

\begin{studybox}{Common Confusions}
\begin{itemize}
  \item A process owns resources and an address space; a thread is scheduled; a fiber is scheduled cooperatively within a thread.
  \item A handle is not a kernel pointer; it is a validated per-process reference with associated access.
  \item An access token describes the caller, while a security descriptor describes the protected object.
  \item The Win32 API is broader than the native system-call interface and can include user-mode work.
\end{itemize}
\end{studybox}

\section*{Key Terms}
\small
\begin{longtable}{C N}
\rowcolor{navy}\color{white}\textbf{Term} & \color{white}\textbf{Concise definition}\\
\endfirsthead
\rowcolor{navy}\color{white}\textbf{Term} & \color{white}\textbf{Concise definition (continued)}\\
\endhead
\textbf{HAL} & Hardware abstraction layer hiding selected platform differences from higher kernel code.\\\hline
\textbf{Executive} & Kernel-mode managers implementing major NT operating-system services.\\\hline
\textbf{Object manager} & Executive component managing typed objects, names, handles, and lifetimes.\\\hline
\textbf{Handle} & Per-process validated reference to an object with granted access.\\\hline
\textbf{Job} & Object grouping processes for accounting and policy.\\\hline
\textbf{Fiber} & Application-scheduled execution context within a thread.\\\hline
\textbf{Section} & Memory object representing mapped files, images, or shared memory.\\\hline
\textbf{Working set} & Process pages currently resident and eligible under residency policy.\\\hline
\textbf{IRP} & Packet describing an I/O operation as it moves through drivers.\\\hline
\textbf{Device stack} & Ordered driver and device objects participating in one logical I/O path.\\\hline
\textbf{MFT} & NTFS table containing one base record for each file or directory.\\\hline
\textbf{Resident attribute} & NTFS attribute stored within an MFT record.\\\hline
\textbf{Access token} & Kernel object holding a subject's identities, groups, privileges, and restrictions.\\\hline
\textbf{Security descriptor} & Object metadata containing ownership, access rules, and audit policy.\\\hline
\textbf{Integrity level} & Mandatory label constraining writes from lower-trust subjects.\\\hline
\textbf{Completion port} & Scalable queue for asynchronous I/O completion notifications.\\\hline
\end{longtable}
\normalsize

\enlargethispage{2\baselineskip}
\begin{studybox}{Active-Recall Review}
\small
\begin{enumerate}
  \item What architectural break separates DOS-based Windows from the NT line?
  \item How do native services and Win32 procedures differ?
  \item What responsibilities belong to the HAL, kernel, and executive?
  \item How does a handle protect kernel object references?
  \item How do jobs, processes, threads, and fibers differ?
  \item What is the relationship among sections, mapped files, and caching?
  \item How does an IRP move through a device stack?
  \item How are resident and nonresident NTFS attributes stored?
  \item How does the reference monitor combine a token and security descriptor?
  \item Why do exploit mitigations complement rather than replace defect removal?
\end{enumerate}
\end{studybox}

\par\medskip
{\color{navy}\bfseries Next-Chapter Connections}\par\smallskip
Chapter 12 abstracts lessons from the case studies into principles for interface design, internal structure, implementation, performance, project organization, and future operating-system directions.
\normalsize
\end{document}
